\documentclass[aspectratio=169]{beamer}
\usepackage{tikz}
\usepackage{graphicx}
\usepackage{listings}
\usepackage{xcolor}
\usepackage{hyperref}
\usepackage[spanish,es-noquoting]{babel}
\usepackage[round]{natbib}

\lstset{
  basicstyle=\ttfamily\small,
  keywordstyle=\color{blue},
  stringstyle=\color{red},
  commentstyle=\color{gray},
  showstringspaces=false
}
% Minimal look
\usetheme{Warsaw}
%\setbeamertemplate{headline}
\setbeamertemplate{footline}{
  \leavevmode%
  \hbox{%
    % Left: Section name
    \begin{beamercolorbox}[wd=.33\paperwidth,ht=2.25ex,dp=1ex,left]{author in head/foot}%
      \hspace*{2ex}\insertsection
    \end{beamercolorbox}%
    % Center: Presentation title
    \begin{beamercolorbox}[wd=.34\paperwidth,ht=2.25ex,dp=1ex,center]{title in head/foot}%
      \insertshorttitle
    \end{beamercolorbox}%
    % Right: Slide numbers
    \begin{beamercolorbox}[wd=.33\paperwidth,ht=2.25ex,dp=1ex,right]{date in head/foot}%
      \insertframenumber{} / \inserttotalframenumber\hspace*{2ex}
    \end{beamercolorbox}%
  }%
  \vskip0pt%
}

\usecolortheme{default}
\usefonttheme{serif}
\graphicspath{{../}{./}}
\setbeamertemplate{caption}[numbered]

% Title info
\title[Module 2: Basic Prompt Enginee]{Module 2: Basic Prompt Engineering}
\subtitle{Key Frameworks for EMs}
\author{Ricardo Reyes-Jimenez}
\institute{Capgemini}
\date{\today}



\begin{document}

% Title slide
\begin{frame}
  \titlepage
\end{frame}

% Reset background for normal slides
\usebackgroundtemplate{}



\begin{frame}[fragile]{Session Objective}


\textbf{By the end of this module, the Engagement Manager will be able to:}

1. Stop communicating with MS Copilot as if it were a search engine.
2. Structure robust prompts using the \textbf{5 Pillars Framework}.
3. Apply the \textbf{TAM Framework} (Task, Audience, Method) to draft precise stakeholder communications.

\end{frame}



\begin{frame}[fragile]{1. The Search Engine Fallacy}


\textbf{AI is not Google. It does not retrieve; it reasons (statistically).}

\begin{itemize}
  \item \textbf{Google Search:} ``Capgemini change request template'' (Keywords).
  \item \textbf{MS Copilot:} ``Act as a Capgemini Delivery Executive. Draft a formal Change Request for a client who added 3 new reporting dashboards to the scope. Use a diplomatic but firm tone.'' (Context + Intent).
\end{itemize}
\textbf{Rule of Thumb:} If your prompt is shorter than a Tweet (280 characters), you are likely under-utilizing the AI and risking hallucinations.

\end{frame}



\begin{frame}[fragile]{2. The 5 Pillars of Prompting}


A vague prompt yields a useless document. Every professional prompt requires 5 pillars:

1. \textbf{Role:} Identity of the AI (e.g., Senior Engagement Manager).
2. \textbf{Context:} Objective and concrete material facts (Project status, client mood).
3. \textbf{Instruction:} The imperative verb (Analyze, Extract, Compare).
4. \textbf{Constraints:} Business limits and rules of what NOT to do.
5. \textbf{Format:} Expected output (Table, 2-page memo, bullet points).

\end{frame}



\begin{frame}[fragile]{Exercise 1: The 5 Pillars Makeover}


\textbf{Context:} Your junior PM asked Copilot: \textit{``Write an email asking for a meeting agenda.''} The result was generic and unprofessional.
\textbf{Task:} Rewrite this prompt using the 5 Pillars.

\begin{itemize}
  \item \textbf{Draft Prompt to Try:}
\end{itemize}
\textit{``[ROLE] Act as an IT Project Manager. [CONTEXT] We have a SteerCo meeting tomorrow with the client's CIO, but they haven't sent the agenda topics yet. [INSTRUCTION] Draft a polite email requesting the agenda. [CONSTRAINTS] Do not sound accusatory; keep it under 3 paragraphs. [FORMAT] Output the email with a clear subject line.''}

\end{frame}



\begin{frame}[fragile]{3. The TAM Framework}


For persuasive communication, structuring the \textit{Why} is as important as the \textit{What}.

\begin{itemize}
  \item \textbf{T - Task:} What exactly needs to be done? (e.g., Explain a 2-week delay).
  \item \textbf{A - Audience:} Who is reading this? What is their technical level and emotional state? (e.g., Non-technical CFO, highly frustrated).
  \item \textbf{M - Method:} How should the information be delivered? (e.g., Executive summary followed by a mitigation plan).
\end{itemize}
\end{frame}



\begin{frame}[fragile]{Exercise 2: Using TAM for Difficult Emails}


\textbf{Context:} The offshore development team missed a critical deadline. You need to inform the client's VP of Engineering without damaging the relationship.
\textbf{Task:} Use the TAM framework in MS Copilot to draft the communication.

\begin{itemize}
  \item \textbf{Draft Prompt to Try:}
\end{itemize}
I need to send an email about a missed deadline.
\textbf{Task:} Explain that the offshore team was delayed by 3 days due to unexpected server downtime, but we are allocating weekend resources to catch up.
\textbf{Audience:} The client's VP of Engineering (technical, results-oriented, currently frustrated).
\textbf{Method:} A direct, professional email starting with the bad news, immediately followed by the mitigation plan, using bullet points for clarity.

\end{frame}



\begin{frame}[fragile]{Recap \& Next Session}


\textbf{Today we learned:}
\begin{itemize}
  \item Stop using keywords; start providing context.
  \item The \textbf{5 Pillars} guarantee structured, predictable outputs.
  \item The \textbf{TAM Framework} tailors the message to the stakeholder's reality.
\end{itemize}
\textbf{Next Session:}
\textit{Module 3: Statement of Work (SOW) Analysis \& Review.} We will learn how to feed long documents into Copilot to extract hidden risks, ambiguous clauses, and SLA commitments.
\end{frame}




\begin{frame}[allowframebreaks]{References}
  \bibliographystyle{apalike}
  \bibliography{references}
\end{frame}


\end{document}